%!TeX root=MemoriaTFG.tex

\chapter{Modelos de lenguaje multimodales: comparación extendida}
\label{cap:anexog}

Este anexo amplía la comparación entre modelos de lenguaje
multimodales generalistas presentada en secci\'on~\ref{sec:llm-results} del
capítulo~\ref{cap:resultados}. Contiene el \emph{prompt} clínico
empleado en la evaluación, el análisis detallado de patrones de
error sobre HAM10000, la caracterización de la sensibilidad al
\emph{prompt} en DermLIP v2 (que motiva el \emph{wrapping}
automático del prototipo) y la comparación de costes operativos
agregados.

\section{Modelos evaluados}
\label{sec:anexoh-modelos}

La tabla~\ref{tab:anexoh-modelos} enumera los modelos evaluados y
sus condiciones de acceso. Los modelos accesibles vía API se han
evaluado con la versión disponible en marzo de 2026; los modelos
locales se han ejecutado sobre la DGX Spark con chip Grace Hopper
GB10 en precisión BF16 sin cuantización adicional.

\begin{table}[H]
\centering
\footnotesize
\setlength{\tabcolsep}{4pt}
\renewcommand{\arraystretch}{0.95}
\caption{Modelos de lenguaje multimodales generalistas evaluados.}
\label{tab:anexoh-modelos}
\begin{tabular}{llrl}
\hline
Modelo                     & Proveedor   & Parámetros & Acceso \\
\hline
GPT-4o                     & OpenAI      & No público  & API REST \\
GPT-4o-mini                & OpenAI      & No público  & API REST \\
Gemini~2.5~Pro             & Google      & No público  & API REST \\
Gemini~2.5~Flash           & Google      & No público  & API REST \\
MedGemma~4B Vision         & Google      & $\sim 4$\,B & Local (Ollama) \\
MedGemma~27B Text          & Google      & $\sim 27$\,B & Local (DGX Spark, BF16) \\
MedGemma~27B + LoRA        & Google      & $\sim 27$\,B + $41$\,M LoRA & Local (DGX Spark, BF16) \\
BLIP-2 Flan-T5-XL          & Salesforce  & ---         & Local (HuggingFace) \\
InstructBLIP Flan-T5-XL    & Salesforce  & ---         & Local (HuggingFace) \\
\hline
\end{tabular}
\end{table}

\section{Prompt clínico estandarizado}
\label{sec:anexoh-prompt}

El \emph{prompt} clínico empleado en la evaluación tiene
aproximadamente $300$ palabras y describe la tarea, presenta
las clases candidatas y solicita el diagnóstico estructurado. La
formulación se mantiene constante entre los nueve proveedores
evaluados, lo que elimina la formulación del \emph{prompt} como
variable confusoria. La estructura es la siguiente:

\begin{quote}
\textit{You are assisting in a dermatology research study. This is
a dermoscopic image of a skin lesion taken during clinical practice.
Your task is to classify the lesion as exactly one of the following
seven categories used in the HAM10000 dataset: melanoma (mel),
melanocytic nevus (nv), basal cell carcinoma (bcc), actinic
keratosis (akiec), seborrheic keratosis (bkl), dermatofibroma (df)
or vascular lesion (vasc). Provide your answer as a single
diagnostic label followed by a brief clinical justification (1--2
sentences). The justification should mention the morphological
features that support the diagnosis (e.g., asymmetry, border
regularity, color pattern, vascular structures). Avoid hedging
language; commit to the most likely diagnosis based on the visible
features. Do not refuse to answer. Respond only with the category
name and the justification.}
\end{quote}

Para Gemini~2.5~Pro y Gemini~2.5~Flash, el \emph{prompt} se ajusta
levemente para evitar la activación de filtros de seguridad
(refusal rate documentado: $\sim 22\%$ en Gemini~2.5~Pro sobre
HAM10000 frente a $<1\%$ en GPT-4o). La etiqueta predicha se
extrae mediante \emph{parsing} determinista que mapea sinónimos
admitidos al código canónico de HAM10000.

\section{Comparativa completa sobre HAM10000, PAD-UFES-20 y DDI}
\label{sec:anexoh-comparativa}

La tabla~\ref{tab:anexoh-completa} reproduce la comparativa de la
tabla~\ref{tab:llm-comparativa} del cuerpo, ampliada con las
variantes Gemini Flash y BLIP-2 + InstructBLIP para incluir el
espectro completo de modelos evaluados.

\begin{table}[H]
\centering
\footnotesize
\setlength{\tabcolsep}{3pt}
\renewcommand{\arraystretch}{0.95}
\caption{Comparativa completa de modelos de lenguaje multimodales
sobre HAM10000 ($N = 1\,232$), PAD-UFES-20 ($N = 461$) y
DDI ($N = 137$).}
\label{tab:anexoh-completa}
\begin{tabular}{llrrrrrr}
\hline
Modelo                     & Modo  & HAM Acc       & HAM BAcc      & PAD Acc       & PAD BAcc      & DDI Acc       & DDI BAcc \\
\hline
MedGemma~27B + LoRA        & FT    & $0{,}802$     & $0{,}270$     & $0{,}430$     & $0{,}347$     & $0{,}657$     & $0{,}618$ \\
MedGemma~27B               & ZS    & $0{,}665$     & $0{,}243$     & $0{,}447$     & $0{,}342$     & $0{,}474$     & $0{,}566$ \\
GPT-4o                     & ZS    & $0{,}485$     & $0{,}465$     & $0{,}553$     & $0{,}523$     & $0{,}723$     & $0{,}648$ \\
Gemini~2.5~Pro             & ZS    & $0{,}403$     & ---           & $0{,}477$     & ---           & $0{,}441$     & --- \\
Gemini~2.5~Flash           & ZS    & $0{,}380$     & ---           & $0{,}486$     & ---           & $0{,}555$     & --- \\
GPT-4o-mini                & ZS    & $0{,}256$     & $0{,}319$     & $0{,}390$     & $0{,}435$     & $0{,}737$     & $0{,}556$ \\
BLIP-2 (Flan-T5-XL)        & ZS    & $0{,}015$     & $0{,}145$     & $0{,}017$     & $0{,}167$     & $0{,}796$     & $0{,}542$ \\
InstructBLIP (Flan-T5-XL)  & ZS    & $0{,}024$     & $0{,}123$     & $0{,}171$     & $0{,}224$     & $0{,}788$     & $0{,}500$ \\
\hline
\end{tabular}
\end{table}

\section{Latencias, refusal rates y costes operativos}
\label{sec:anexoh-coste}

La tabla~\ref{tab:anexoh-coste} reporta las características
operativas de la evaluación: latencia media por imagen, tasa de
respuestas no clasificables (\emph{refusal rate}) y coste agregado
sobre los tres datasets evaluados.

\begin{table}[H]
\centering
\footnotesize
\setlength{\tabcolsep}{3pt}
\renewcommand{\arraystretch}{0.95}
\caption{Características operativas de la evaluación de LLMs
multimodales. \emph{Refusal} sobre Gemini~2.5~Pro corresponde a
respuestas bloqueadas por filtros de seguridad. ``Coste'' agrega
HAM10000 + PAD-UFES-20 + DDI.}
\label{tab:anexoh-coste}
\begin{tabular}{lrrrrr}
\hline
Modelo               & Latencia/imagen & Refusal HAM & Refusal PAD & Refusal DDI & Coste total \\
\hline
GPT-4o               & $\sim 4$\,s     & $<1\%$       & $<1\%$       & $<1\%$       & $\sim$\$2{,}71 \\
GPT-4o-mini          & $\sim 3$\,s     & $<1\%$       & $<1\%$       & $<1\%$       & $\sim$\$1{,}36 \\
Gemini~2.5~Pro       & $\sim 8$\,s     & $21{,}5\%$   & $15{,}0\%$   & $18{,}2\%$   & --- \\
Gemini~2.5~Flash     & $\sim 2$\,s     & $0{,}3\%$    & $2{,}4\%$    & $2{,}7\%$    & --- \\
MedGemma~27B         & $1{,}76$\,s     & $0\%$        & $0\%$        & $0\%$        & $0$\,EUR \\
MedGemma~27B + LoRA  & $1{,}82$\,s     & $0\%$        & $0\%$        & $0\%$        & $0$\,EUR (entrenam.) \\
\hline
\end{tabular}
\end{table}

Tres observaciones operativas: (i)~Gemini~2.5~Pro presenta
\emph{refusal rate} elevado sobre HAM10000 ($21{,}5\%$),
comportamiento atribuible a la activación de filtros de seguridad
sobre imagen clínica. (ii)~MedGemma~27B y su variante con LoRA
operan localmente sin coste recurrente, con latencias medidas
sobre la DGX Spark de $\sim 1{,}8$ segundos por imagen,
inferiores a las de los proveedores comerciales. (iii)~El coste
agregado de GPT-4o-mini ($\sim \$1{,}36$ sobre $1\,830$
imágenes) es de aproximadamente $\$0{,}001$ por imagen, lo que
hace operativamente viable la evaluación a escala
moderada; el coste de GPT-4o ($\sim \$2{,}71$) duplica esta
cifra.

\section{Patrones de error sobre HAM10000}
\label{sec:anexoh-errores}

La tabla~\ref{tab:anexoh-confusion-gpt4o} reporta la matriz de
predicciones de GPT-4o sobre HAM10000. La fila ``\emph{Predicciones
totales}'' indica el número de imágenes asignadas a cada clase por
el modelo; la comparación con la fila ``$N_{\mathrm{test}}$''
revela los sesgos sistemáticos.

\begin{table}[H]
\centering
\footnotesize
\setlength{\tabcolsep}{4pt}
\renewcommand{\arraystretch}{0.95}
\caption{Distribución de predicciones de GPT-4o sobre HAM10000.}
\label{tab:anexoh-confusion-gpt4o}
\begin{tabular}{lrrr}
\hline
Clase                       & $N_{\mathrm{test}}$ & Predicciones GPT-4o & Factor \\
\hline
actinic keratosis           & $35$    & $2$    & $0{,}06\times$ \\
basal cell carcinoma        & $44$    & ---    & --- \\
seborrheic keratosis        & $107$   & ---    & --- \\
dermatofibroma              & $8$     & $284$  & $35{,}5\times$ \\
melanoma                    & $70$    & $299$  & $4{,}3\times$ \\
melanocytic nevus           & $951$   & ---    & --- \\
vascular lesion             & $17$    & ---    & --- \\
\hline
\end{tabular}
\end{table}

El sesgo más acusado es la sobreasignación de \emph{dermatofibroma}
(predicciones $284$ frente a $8$ muestras reales, factor
$35{,}5\times$). La hipótesis sobre el mecanismo es que el modelo
responde al espacio léxico del \emph{prompt}: la enumeración
explícita de \emph{dermatofibroma} entre las clases candidatas
activa preferentemente esta categoría como respuesta cuando la
evidencia visual no es saliente. Un patrón análogo, aunque de
menor magnitud, se observa con melanoma (factor $4{,}3\times$),
clase clínicamente prioritaria que el modelo podría sobreasignar
por reflejar un sesgo precautorio aprendido durante el
entrenamiento.

GPT-4o-mini presenta un patrón distinto: sesgo masivo hacia
\emph{seborrheic keratosis} ($528$ predicciones frente a
$107$ muestras reales) y melanoma ($359$ predicciones frente
a $70$). Sólo $196$ de los $951$ casos de
\emph{melanocytic nevus} se clasifican correctamente
(\emph{recall} $0{,}196$), mientras GPT-4o alcanza
\emph{recall} de $0{,}505$ sobre la misma clase.

\section{MedGemma~27B con LoRA: análisis por dataset}
\label{sec:anexoh-medgemma}

La adaptación con LoRA sobre MedGemma~27B
(\emph{rank} $r = 16$, $\alpha = 32$, dropout $0{,}05$,
$1$ época, $\eta = 2\times 10^{-5}$, \emph{effective batch}
$8$, $\sim 41$\,M parámetros entrenables sobre el
\emph{split} de entrenamiento de HAM10000) produce las mejoras
asimétricas reportadas en secci\'on~\ref{sec:disc-llm}. La
tabla~\ref{tab:anexoh-medgemma-perclass} reporta el desglose por
clase sobre HAM10000 con la variante adaptada.

\begin{table}[H]
\centering
\footnotesize
\setlength{\tabcolsep}{4pt}
\renewcommand{\arraystretch}{0.95}
\caption{Desglose por clase de MedGemma~27B + LoRA sobre HAM10000
(\emph{precision}, \emph{recall}, F1).}
\label{tab:anexoh-medgemma-perclass}
\begin{tabular}{lrrrr}
\hline
Clase                       & $N_{\mathrm{test}}$ & Precision & Recall    & F1 \\
\hline
actinic keratosis           & $35$    & $0{,}32$  & $0{,}26$  & $0{,}29$ \\
basal cell carcinoma        & $44$    & $0{,}34$  & $0{,}25$  & $0{,}29$ \\
dermatofibroma              & $8$     & $0{,}00$  & $0{,}00$  & $0{,}00$ \\
melanocytic nevus           & $951$   & $0{,}85$  & $0{,}99$  & $0{,}92$ \\
melanoma                    & $70$    & $0{,}38$  & $0{,}36$  & $0{,}37$ \\
seborrheic keratosis        & $107$   & $0{,}67$  & $0{,}04$  & $0{,}07$ \\
vascular lesion             & $17$    & $0{,}00$  & $0{,}00$  & $0{,}00$ \\
\hline
\end{tabular}
\end{table}

La concentración del rendimiento sobre la clase mayoritaria
(nevus melanocítico: \emph{recall} $0{,}99$) a costa de las
minoritarias (dermatofibroma y vascular lesion: F1 $0{,}00$)
contrasta con el comportamiento de PanDerm Base con ajuste fino y TTA
sobre el mismo dataset, que mantiene \emph{recall} superior a
$0{,}65$ sobre todas las clases minoritarias gracias al
\emph{weighted random sampler} de la receta original
(secci\'on~\ref{sec:ft}). La diferencia metodológica sugiere que el
\emph{fine-tuning} de LLMs multimodales sobre datasets
desbalanceados requiere estrategias específicas de balanceo no
contempladas en la configuración estándar de LoRA.

\section{Sensibilidad al prompt en DermLIP v2}
\label{sec:anexoh-prompt-derm}

El módulo M6 del prototipo DermApIxel
(secci\'on~\ref{sec:dermapixel-m6}) opera sobre DermLIP v2 con
\emph{wrapping} automático del \emph{prompt}. Esta sección
documenta cuantitativamente el fenómeno que motiva el
\emph{wrapping}.

\begin{table}[H]
\centering
\footnotesize
\setlength{\tabcolsep}{4pt}
\renewcommand{\arraystretch}{0.95}
\caption{Sensibilidad al \emph{prompt} de DermLIP v2 sobre
HAM10000: curva de AUROC observada en función de la longitud
del \emph{prompt}.}
\label{tab:anexoh-prompt-curva}
\begin{tabular}{lll}
\hline
Longitud del \emph{prompt}              & Estrategia                                         & AUROC HAM10000 \\
\hline
$1$ palabra                             & Etiqueta desnuda                                    & $\approx$ aleatorio \\
$\sim 5$ palabras                       & Envoltorio clínico ($``dermoscopy image of \{X\}''$) & $\sim 0{,}72$ \\
$8$--$15$ palabras                      & Envoltorio + $1$ rasgo clínico                      & $\mathbf{0{,}854}$ \\
$\sim 20$ palabras                      & Envoltorio + $3$ conceptos Derm1M                   & $0{,}842$ \\
$> 25$ palabras                         & Envoltorio + $5$ conceptos Derm1M                   & $0{,}826$ \\
\hline
\end{tabular}
\end{table}

La curva presenta un punto óptimo en $8$--$15$ palabras con
un único rasgo clínico añadido (AUROC $0{,}854$). El
rendimiento decae tanto en el régimen de \emph{prompt} muy corto
(etiqueta desnuda: AUROC aleatorio) como en el régimen de
\emph{prompt} muy largo ($>25$ palabras: AUROC $0{,}826$).

La causa propuesta sobre el régimen de \emph{prompt} muy corto es
que PubMedBERT (encoder textual de DermLIP v2) está preentrenado
sobre \emph{captions} del corpus Derm1M con estructura
``dermoscopy image of $[\text{clase}]$, $[\text{descripción}]$''.
Una palabra aislada se proyecta sobre regiones del espacio de
\emph{embedding} muy poco pobladas, lo que produce similitudes
coseno arbitrarias frente al \emph{embedding} visual. Sobre el
régimen de \emph{prompt} muy largo, la hipótesis es que la
acumulación de conceptos clínicos introduce ruido semántico que
diluye la señal de la clase principal.

El \emph{wrapping} automático del prototipo implementa esta curva
en tres capas: presets predefinidos con \emph{prompts} ya en el
punto óptimo, composable \texttt{useZeroShot.js} que aplica el
envoltorio $``dermoscopy image of \{term\}''$ cuando la entrada
del usuario contiene cinco o menos palabras, y conmutador en la
interfaz que permite al dermatólogo desactivar el envoltorio si
escribe descripciones completas propias. La caracterización
\textbf{[NO ESPECIFICADO: ampliar el barrido sistemático de
longitud del prompt sobre los demás datasets de evaluación]}
constituye una línea de continuación inmediata.

\section{Síntesis}
\label{sec:anexoh-sintesis}

La evaluación extendida documentada en este anexo confirma los
patrones cualitativos identificados en secci\'on~\ref{sec:disc-llm} del
capítulo~\ref{cap:discusion}: brecha de $43{,}5$\,pp accuracy
entre el mejor especialista ajustado (PanDerm Base FT con TTA) y el
mejor LLM multimodal generalista \emph{zero-shot} (GPT-4o); sesgo
léxico sistemático en los modelos cerrados; \emph{refusal rates}
elevados en Gemini~2.5~Pro que limitan su utilidad clínica; y
asimetría en la adaptación con LoRA sobre MedGemma~27B. Las
cifras detalladas refuerzan la conclusión de que los LLMs
multimodales generalistas no constituyen sustitutos defendibles
de los encoders fundacionales específicos del dominio en el
régimen clínico evaluado, pero pueden complementarlos como
generadores de razonamiento textual estructurado bajo el
\emph{prompt} clínico descrito en secci\'on~\ref{sec:anexoh-prompt}.
